\section{Related Work}

There has been growing interest in developing resources for Qur’anic text and recitation, yet existing efforts typically focus on either textual analysis or audio recordings, rather than providing a fully integrated multimodal resource. Early structured corpora, such as the \emph{Quranic Arabic Corpus} \cite{dukhovny2009qac}, provide full text with detailed morphological and syntactic annotations, supporting parsing, part-of-speech tagging, and grammatical analysis. More recent contributions, including the \emph{MASAQ corpus} \cite{sawalha2024masaq} and the corpus introduced by \cite{nashir2025quranic}, expand these annotations and incorporate orthographic variants (Uthmani and Imlaai scripts), Buckwalter and phonetic transliterations, as well as English translations at the verse level. Similarly, the \emph{Tanzil project} \cite{tanzil2008} offers the Qur’an aligned with translations in more than 40 languages, supporting multilingual NLP research, though without detailed linguistic or recitational information.

\begin{table*}[t]
\centering
\caption{Feature comparison across prominent Quranic datasets. The proposed \textsc{Qur’an-MD} is the first to holistically integrate all listed modalities across multiple reciters at both verse and word levels, filling a critical gap in the available resources.}
\label{tab:comparison_final}
\resizebox{0.9\textwidth}{!}{%
\begin{tabular}{@{}lccccc@{}}
\toprule
 & \multicolumn{5}{c}{\textbf{Dataset}} \\
\cmidrule(l){2-6}
\textbf{Feature} & \begin{tabular}[c]{@{}c@{}}Quranic Arabic \\ Corpus \cite{dukhovny2009qac}\end{tabular} & \begin{tabular}[c]{@{}c@{}}Tanzil \\ Project \cite{tanzil2008}\end{tabular} & Tarteel \cite{khan2021tarteel} & QDAT \cite{qdat2021} & \begin{tabular}[c]{@{}c@{}}\textbf{\textsc{Qur’an-MD}} \\ \textbf{(Ours)}\end{tabular} \\ 
\midrule
Text (Arabic) & \yes & \yes & \yes & \yes & \yes \\
Morpho-syntactic Data & \yes & \no & \no & \no & \yes \\
Multi-Language Translation & \no & \yes & \no & \no & \yes \\
Word-level Audio & \no & \no & \no & \yes & \yes \\
Verse-level Audio & \no & \no & \yes & \yes & \yes \\
Multiple Reciters & \no & \no & \yes & \no & \yes \\ 
\bottomrule
\end{tabular}%
}
\end{table*}


Parallel to these textual resources, several initiatives have sought to capture Qur’anic recitation across multiple readers. Large-scale audio collections, such as \emph{Tarteel AI EveryAyah}\cite{khan2021tarteel} and the \emph{Quran-Recitations Dataset} \cite{quran_recitations_hf_dataset}, provide verse-level recordings across numerous reciters, while the \emph{Quranic Audio Dataset} \cite{quran_audio_dataset2024} and \emph{RetaSy} \cite{salameh2024retasy} focus on non-native reciters with labeled correctness annotations. The \emph{OpenSLR Quran Speech-to-Text corpus} \cite{openslr132} offers complete verse coverage for multiple reciters, and smaller targeted collections, such as \emph{QDAT} \cite{osman2021qdat} and its extensions \cite{eval_pronunciation_tajweed2025}, provide labeled data for Tajweed error detection using deep neural models. Recent Kaggle contributions, including the \emph{Quran Ayat Speech-to-Text} dataset \cite{bigguyubuntu2023ayat}, \emph{Quran Reciters dataset} \cite{omartariq2025}, \emph{Quran.com Audio dataset} \cite{abdo3id2022qurancom}, \emph{Quran Recitations for Audio Classification dataset} \cite{alrajeh2021recitations}, and the \emph{Comprehensive Quranic Dataset v1 (CQDV1)} \cite{cqdv12024}, offer large-scale verse-level audio recordings from dozens of reciters, but generally provide only Arabic text and audio without consistent transliteration or translation.

Taken together, these resources demonstrate substantial progress in both linguistic annotation and audio collection, yet they remain fragmented: either focusing on text or audio, verse- or word-level data. None combine verse and word-level text, transliteration, English translation, and audio across multiple reciters with consistent cross-modal alignment. \textsc{Qur’an-MD} addresses these gaps by integrating verse and word-level text, English translation, transliteration, word-level audio, and verse audio from 30 reciters, providing a unified multimodal resource that supports research across natural language processing, speech technologies, linguistic analysis, and digital Islamic studies.
